\newpage
\chapter*{Rezumat}
\addcontentsline{toc}{chapter}{Rezumat}

Această lucrare prezintă proiectarea și implementarea unei platforme web pentru managementul fluxurilor studențești și monitorizarea prezenței la activitățile de laborator. Punctul de plecare a fost experiența practică din mediul universitar, unde utilizarea platformelor existente acoperă bine zona de e-learning, dar lasă adesea procesele operaționale distribuite între mai multe instrumente și evidențe separate.

Soluția propusă, \textbf{Student Management System}, urmărește să unifice aceste fluxuri într-o aplicație coerentă, cu arhitectură client-server: frontend Single Page Application dezvoltat în Angular, backend REST implementat în PHP și persistență de date în MySQL. Din perspectiva tehnică, accentul este pus pe mecanismele Angular care susțin scalabilitatea și mentenabilitatea aplicației: organizare modulară pe componente, servicii și Dependency Injection, routing cu protecția rutelor, validare prin Reactive Forms și comunicare reactivă cu RxJS.

La nivel funcțional, aplicația include autentificare pe roluri, administrarea studenților și a grupelor, gestionarea cursurilor și laboratoarelor, marcarea prezenței pe sesiuni și importul de date din fișiere Excel. Validarea a fost realizată prin scenarii de utilizare care urmăresc fluxurile reale de lucru pentru administratori, profesori și studenți.

Rezultatul este o aplicație modulară, orientată spre utilizare practică, care reduce operațiile repetitive, crește coerența datelor și oferă o bază solidă pentru extinderi ulterioare, precum optimizări de performanță, consolidarea testării automate și integrarea unor funcționalități avansate de raportare.

\textbf{Cuvinte-cheie:} Angular, REST API, PHP, MySQL, management studenți, monitorizare prezență, aplicație web.